\documentclass[a4paper]{arrowhead}

\usepackage[yyyymmdd]{datetime}
\usepackage{etoolbox}
\usepackage[utf8]{inputenc}
\usepackage[T1]{fontenc}
\usepackage{multirow}
\usepackage{longtable}
\usepackage{tikz}
\usetikzlibrary{positioning,shapes.multipart,arrows.meta}

\renewcommand{\dateseparator}{-}
\setlength{\emergencystretch}{3em}
\lstset{language=HTTP,breakatwhitespace=false,columns=fullflexible}

%% Special references
\newcommand{\mref}[2][]{{\textcolor{ArrowheadPurple}{\hyperref[sec:model:#2]{\ifstrempty{#1}{#2}{#1}}}}}
\newcommand{\pref}[1]{{\textcolor{ArrowheadGrey}{\hyperref[sec:model:primitives:#1]{#1}}}}
%%

\usepackage{xurl}
\begin{document}
\sloppy

%% Arrowhead Document Properties
\ArrowheadTitle{Choreographer}
\ArrowheadType{System Description}
\ArrowheadTypeShort{SysD}
\ArrowheadVersion{5.3.0}
\ArrowheadDate{2026-11-09}
\ArrowheadAuthor{Jerker Delsing}
\ArrowheadStatus{REVIEW DRAFT}
\ArrowheadContact{jerker.delsing@ltu.se}
\ArrowheadFooter{\href{www.arrowhead.eu}{www.arrowhead.eu}}
\ArrowheadSetup
%%

%% Front Page
\begin{center}
  \vspace*{1cm}
  \huge{\arrowtitle}

  \vspace*{0.2cm}
  \LARGE{\arrowtype}
  \vspace*{1cm}

  \vspace*{\fill}

  % Front Page Image

  \vspace*{1cm}
  \vspace*{\fill}

  % Front Page Abstract
  \begin{abstract}
    This document describes the Choreographer support system of the Eclipse Arrowhead framework (documentation edition 5.3.0), as implemented in \texttt{ah5-Choreographer-java-spring} of the primary Java source tree (\texttt{Eclipse-Arrowhead-java}) and deployed with the standalone Compose definition of that folder.
  \end{abstract}

  \vspace*{1cm}
\end{center}

\newpage
%%

%% Table of Contents
\tableofcontents
\newpage
%%

\section{Documentation edition and implementation scope}
This is the 5.3.0 documentation edition, dated 2026-11-09. It carries forward the system and service contracts from the retained 5.2.2 documentation and identifies the current primary implementation sources. The former \path|documentation/v5.2.2| folder is retained unchanged.

Document version 5.3.0 is distinct from executable, dependency, registered system/service and wire-contract versions. Values still defined by the source (including 5.2.2 artefact versions, 1.0.0 service versions and \path|SysMonitor-AMQP-IDD-5.2.2| payload identifiers) are preserved. This documentation update does not change runtime behaviour or establish a new implementation release. Historical verification statements describe the tests recorded in the inherited baseline; this revision validates document compilation and references.

\section{Overview}
\label{sec:overview}

This document describes the Choreographer system. The Choreographer executes workflow plans: for each step it selects an executor system, has Dynamic\allowbreak{}Service\allowbreak{}Orchestration orchestrate the required provider for the executor, starts the executor and follows the results.

\begin{figure}[ht!]
  \centering
  \begin{tikzpicture}[font=\small, >=Stealth]
    \node[draw=ArrowheadBlue, very thick, minimum width=4.2cm, minimum height=2.2cm, align=center, fill=ArrowheadBlue!10] (s)
      {\guillemotleft{}System\guillemotright{}\\\textbf{Choreographer}};
    \node[draw=ArrowheadTeal, rectangle, align=left, left=1.3cm of s, text width=4.6cm, inner sep=5pt] (p)
      {\textbf{Produced services}\\ \texttt{choreographer\allowbreak{}Session}\\ \texttt{choreographer\allowbreak{}Executor}\\ \texttt{choreographer\allowbreak{}Management}\\ \texttt{generalManagement}\\ \texttt{sysMonitor}};
    \node[draw=ArrowheadOrange, rectangle, align=left, right=1.3cm of s, text width=4.6cm, inner sep=5pt] (c)
      {\textbf{Consumed}\\ AMQP broker (SysMonitor)\\ Authentication\\ ConsumerAuthorization\\ DynamicServiceOrchestration\\ Executors\\ ServiceRegistry};
    \draw[->, thick, ArrowheadTeal] (s) -- (p);
    \draw[->, thick, ArrowheadOrange] (s) -- (c);
  \end{tikzpicture}
  \caption{SysML-style block view of Choreographer with its produced and consumed services.}
  \label{fig:system}
\end{figure}

The rest of this document is organized as follows.
In Section \ref{sec:prior_art}, we reference major prior art capabilities of the system.
In Section \ref{sec:use}, we describe the intended usage of the system.
In Section \ref{sec:properties}, we describe fundamental properties provided by the system.
In Section \ref{sec:delimitations}, we describe delimitations of capabilities of the system.
In Section \ref{sec:services}, we describe the abstract service functions consumed or produced by the system.
In Section \ref{sec:security}, we describe the security capabilities of the system.

\newpage

\subsection{Significant Prior Art}
\label{sec:prior_art}

Successor of the Arrowhead 4.6 Choreographer core system.

\subsection{How This System Is Meant to Be Used}
\label{sec:use}

Operators create plans; clients or operators start sessions; executors register and report step results. DSO pushes orchestration results to \texttt{POST /\allowbreak{}choreographer/\allowbreak{}orchestration-\allowbreak{}result/\allowbreak{}\{correlation\allowbreak{}Id\}} (not a registered service).

\subsection{System functionalities and properties}
\label{sec:properties}

\subsubsection{Functional properties of the system}

\begin{itemize}
  \item Plan validation (cycles, reachability, transitive reduction).
  \item Session engine with parallel independent steps.
  \item Executor selection and retry on \texttt{ERROR}; abort on \texttt{FATAL\_\allowbreak{}ERROR}.
  \item Push orchestration for executors through DSO.
\end{itemize}

\subsubsection{Configuration of system properties}

The system is configured with Spring properties (\texttt{application.properties}, overridden by environment variables in the container deployment). The properties that shape its behaviour are:

\noindent\begin{longtable}{| p{5.4cm} | p{3.4cm} | p{6.4cm} |} \hline
\rowcolor{gray!33} Property & Default & Meaning \\ \hline \endhead
\texttt{server.\allowbreak{}port} & 8457 & HTTPS port. \\ \hline
\texttt{authentication.\allowbreak{}policy} & declared & Standalone deployment: \texttt{outsourced} with Authentication identity tokens and mutual TLS. \\ \hline
\texttt{management.\allowbreak{}policy} / \texttt{management.\allowbreak{}whitelist} & authorization & Access to the management services. \\ \hline
\texttt{enable.\allowbreak{}blacklist.\allowbreak{}filter} / \texttt{force.\allowbreak{}blacklist.\allowbreak{}filter} & false / true & Blacklist check of requesters. \\ \hline
\texttt{max.\allowbreak{}plan.\allowbreak{}iteration} & 1000 & Largest session quantity. \\ \hline
\texttt{orchestration.\allowbreak{}timeout} & 30 & Wait for a pushed orchestration result (s). \\ \hline
\texttt{session.\allowbreak{}worker.\allowbreak{}threads} & 8 & Engine worker threads. \\ \hline
\texttt{normalization.\allowbreak{}mode} & extended & Name normalization. \\ \hline
\texttt{sysmonitor.\allowbreak{}enabled,\allowbreak{} sysmonitor.\allowbreak{}broker.\allowbreak{}*,\allowbreak{} sysmonitor.\allowbreak{}ssl.\allowbreak{}*} & true, amqp-broker:5671, server.ssl stores & Native \texttt{sysMonitor} publisher of the common library. \\ \hline
\texttt{arrowhead.\allowbreak{}ssl.\allowbreak{}crl-\allowbreak{}file} & (unset) & CRL checked for every client certificate (standalone deployment: \texttt{/\allowbreak{}crl/\allowbreak{}ca.\allowbreak{}crl}). \\ \hline
\end{longtable}

Secrets (database, keystore and identity passwords, secret keys) have no usable defaults; the standalone deployment reads them from per-deployment Compose file secrets under \texttt{/\allowbreak{}run/\allowbreak{}secrets}.

The standalone deployment of this folder uses \texttt{config/\allowbreak{}application.\allowbreak{}properties} (copied into the image) with \texttt{server.\allowbreak{}port}=8454 (source default 8457), \texttt{authentication.\allowbreak{}policy}=\texttt{outsourced}, \texttt{service.\allowbreak{}interface.\allowbreak{}policy}=\texttt{CERT\_AUTH}, \texttt{management.\allowbreak{}policy}=\texttt{sysop-\allowbreak{}only}, \texttt{service.\allowbreak{}authorization.\allowbreak{}enabled}=\texttt{true}, \texttt{mqtt.\allowbreak{}api.\allowbreak{}enabled}=\texttt{false}, \texttt{arrowhead.\allowbreak{}ssl.\allowbreak{}crl-\allowbreak{}file}=\texttt{/\allowbreak{}crl/\allowbreak{}ca.\allowbreak{}crl}; secrets are read from Compose file secrets under \texttt{/\allowbreak{}run/\allowbreak{}secrets} (\texttt{spring.\allowbreak{}config.\allowbreak{}import=\allowbreak{}configtree:\allowbreak{}/\allowbreak{}run/\allowbreak{}secrets/}). \texttt{enable.\allowbreak{}blacklist.\allowbreak{}filter} is not enabled there, so incoming requests are not checked against the Blacklist; blacklisted providers are still removed by Dynamic\allowbreak{}Service\allowbreak{}Orchestration.

\subsubsection{Data stored by the system}

Plans, actions, steps, executors, sessions, session steps, worklog and log entries (database \texttt{ah\_\allowbreak{}choreographer}).

\subsubsection{Non functional properties}

Steps run on a bounded worker pool; orchestration results time out after \texttt{orchestration.\allowbreak{}timeout}.

\subsubsection{Stateful or stateless}

Stateful; sessions active at restart are aborted.

\subsection{Important Delimitations}
\label{sec:delimitations}

No inter-cloud executor discovery; inter-cloud only through the orchestration flag.

\newpage

\section{Services}
\label{sec:services}

This section describes the produced and consumed services. Each produced service is specified by an SD document and one IDD document per interface.

\subsection{Produced services}

\noindent\begin{longtable}{| p{5.3cm} | p{1.3cm} | p{2.6cm} | p{1.6cm} | p{3.6cm} |} \hline
\rowcolor{gray!33} Service & Version & Interfaces & SD & IDD \\ \hline \endhead
\texttt{choreographer\allowbreak{}Session} & 1.0.0 & HTTP(S) & \cite{cho-choreographerSession-sd} & \cite{cho-choreographerSession-http-idd} \\ \hline
\texttt{choreographer\allowbreak{}Executor} & 1.0.0 & HTTP(S) & \cite{cho-choreographerExecutor-sd} & \cite{cho-choreographerExecutor-http-idd} \\ \hline
\texttt{choreographer\allowbreak{}Management} & 1.0.0 & HTTP(S) & \cite{cho-choreographerManagement-sd} & \cite{cho-choreographerManagement-http-idd} \\ \hline
\texttt{generalManagement} & 1.0.0 & HTTP(S) & \cite{cho-generalManagement-sd} & \cite{cho-generalManagement-http-idd} \\ \hline
\texttt{sysMonitor} & 1.0.0 & AMQP(S) & \cite{cho-sysMonitor-sd} & \cite{cho-sysMonitor-amqp-idd} \\ \hline
\end{longtable}

The HTTP interfaces are always offered and the AMQP \texttt{sysMonitor} interface with \texttt{sysmonitor.enabled=true}. The unregistered endpoints of the implementation are: \texttt{POST /\allowbreak{}choreographer/\allowbreak{}orchestration-\allowbreak{}result/\allowbreak{}\{correlation\allowbreak{}Id\}}.

\subsection{Consumed services}

\noindent\begin{longtable}{| >{\raggedright\arraybackslash}p{2.9cm} | >{\raggedright\arraybackslash}p{3.1cm} | >{\raggedright\arraybackslash}p{2.5cm} | p{3.6cm} | >{\raggedright\arraybackslash}p{2.9cm} |} \hline
\rowcolor{gray!33} Provider & Service & Operations & Use & Condition \\ \hline \endhead
ServiceRegistry & systemDiscovery, serviceDiscovery & register, revoke, lookup & Registration of the system and its services; lookup of the bootstrap services. & always \\ \hline
Authentication & identity & identity-login, identity-verify & Own identity token; verification of requester tokens. & \texttt{authentication.\allowbreak{}policy=\allowbreak{}outsourced} \\ \hline
Consumer\allowbreak{}Authorization & authorization & verify & Management filter decisions. & \texttt{management.\allowbreak{}policy=\allowbreak{}authorization} \\ \hline
Dynamic\allowbreak{}Service\allowbreak{}Orchestration & service\allowbreak{}Orchestration\allowbreak{}Push\allowbreak{}Management & subscribe, trigger & Orchestration for executors. & grant \\ \hline
Executors & executor interface & service-info, start, abort & Step execution. & sessions \\ \hline
AMQP broker (SysMonitor) & sysMonitor (publisher) & echo, power, status & Own observations. & \texttt{sysmonitor.\allowbreak{}enabled=\allowbreak{}true} \\ \hline
\end{longtable}

Except for the bootstrap services, which are looked up in the ServiceRegistry, the system reaches other systems' services through DynamicServiceOrchestration (\texttt{http.collector.mode=ORCHESTRATION}, results cached for 300 s); every such call needs a ConsumerAuthorization grant.

\newpage

\section{Security}
\label{sec:security}

The system can run in an insecure mode (\texttt{server.ssl.enabled=false}, plain HTTP, identity declared by the requester) for development, and in the Arrowhead secure mode with HTTPS and X.509 certificates following the Arrowhead certificate profile (\texttt{<System>.<Cloud>.<Organization>.arrowhead.eu}). Certificates not issued under the cloud CA are refused by the TLS handshake. The standalone deployment of this folder (\texttt{config/\allowbreak{}application.\allowbreak{}properties} and \texttt{compose.yaml}) runs the secure mode.

\subsection{Security model}


\begin{itemize}
  \item \textbf{Protocols:} HTTPS (TLS 1.3 and 1.2) with mandatory client certificates; MQTT over TLS; AMQP over TLS for \texttt{sysMonitor}.
  \item \textbf{Data protection:} transport encryption by TLS; the database connection uses TLS with server verification in the standalone deployment. Secrets are kept out of the container environment (Compose file secrets).
  \item \textbf{System authentication:} by the authentication policy: \texttt{certificate}, \texttt{outsourced} (identity token verified with Authentication) or \texttt{declared} (insecure); the standalone deployment uses \texttt{outsourced}. An identity token is accepted only together with the client certificate of the system it was issued to (token-to-certificate binding), and every client certificate is checked against the cloud CA's CRL (fail closed).
  \item \textbf{Authorization of produced services:} management services are protected by the management filter (\texttt{management.policy}: \texttt{SYSOP\_ONLY}, \texttt{WHITELIST} or \texttt{AUTHORIZATION}); with \texttt{service.authorization.enabled=true} ConsumerAuthorization decides each call; an optional blacklist filter refuses blacklisted requesters. Sysop identities bypass the management filter.
  \item The orchestration-result endpoint accepts only Dynamic\allowbreak{}Service\allowbreak{}Orchestration's certificate and pending random correlation ids.
\end{itemize}

For the Arrowhead certificate profile see github.com/eclipse-arrowhead/documentation.


\bibliographystyle{IEEEtran}
\bibliography{bibliography}

\newpage

\section{Revision History}
\subsection{Amendments}

\noindent\begin{tabularx}{\textwidth}{| p{1cm} | p{3cm} | p{2cm} | X | p{4cm} |} \hline
\rowcolor{gray!33} No. & Date & Version & Subject of Amendments & Author \\ \hline

1 & 2026-10-22 & 5.2.2 & Document created from the Arrowhead LaTeX templates for the 5.2.2 Java Spring implementation, including the local extensions of this workspace. & Jerker Delsing \\ \hline
2 & 2026-10-30 & 5.2.2 & Checked against the primary Java source tree (Eclipse-Arrowhead-java): standalone deployment, configuration and security statements updated. & Claude Opus 5.5 (draft) \\ \hline
3 & 2026-11-09 & 5.3.0 & New documentation edition; primary-source references and Go/Rust implementation scope; 5.2.2 retained. & Documentation maintenance \\ \hline
\end{tabularx}

\subsection{Quality Assurance}

\noindent\begin{tabularx}{\textwidth}{| p{1cm} | p{3cm} | p{2cm} | X |} \hline
\rowcolor{gray!33} No. & Date & Version & Approved by \\ \hline

1 & & \arrowversion  &  \\ \hline

\end{tabularx}

\end{document}
